\section{Space Complexity --- Machines --- Program}
%%% generated by the blueprint pipeline from TCSlib.Complexity.SpaceComplexity.Machines.Program
\subsection*{Space Complexity --- Machines --- Program: Introduction}

Declarations of the Lean module \texttt{TCSlib.Complexity.SpaceComplexity.Machines.Program}.

\subsection*{Space Complexity --- Machines --- Program: Results}

\begin{definition}[Virtual input mode]
\lean{Complexity.LogProg.Mode}
\label{Complexity.LogProg.Mode}
\leanok
A virtual-input mode is one of two choices: use the whole input, doubled, or use the
leading run of $1$s in the input as a plain segment.
\end{definition}

\begin{definition}[First virtual-input segment]
\lean{Complexity.LogProg.Mode.seg0}
\label{Complexity.LogProg.Mode.seg0}
\leanok
\uses{Complexity.LogProg.Mode, Complexity.LogProg.Seg}
For a mode and Boolean word $x$, the first segment is the pair $(x,1)$ in whole-input
mode, and the pair consisting of the leading run of $1$s in $x$ and $0$ in unary-prefix
mode.
\end{definition}

\begin{definition}[Argument word segmentation]
\lean{Complexity.LogProg.argSegs}
\label{Complexity.LogProg.argSegs}
\leanok
\uses{Complexity.LogProg.Seg}
For a list of Boolean words, the argument segments are empty for an empty list, and
consist of the sole word paired with $0$ for a one-word list. For a list with at least
two words, the first word is paired with $1$, followed by the argument segments of the
remaining words.
\end{definition}

\begin{lemma}[Length of the Argument Segmentation]
\lean{Complexity.LogProg.length_argSegs}
\label{Complexity.LogProg.length_argSegs}
\leanok
\uses{Complexity.LogProg.argSegs}
\difficulty{4}
For every finite list of Boolean words, the segmentation into argument segments has the
same length as the original list.
\end{lemma}

\begin{lemma}[Segments formed from argument words]
\lean{Complexity.LogProg.seg_argSegs}
\label{Complexity.LogProg.seg_argSegs}
\leanok
\uses{Complexity.LogProg.argSegs, Complexity.LogProg.seg, Complexity.LogProg.seg_cons_succ}
\difficulty{5}
Let $w_0,\ldots,w_{n-1}$ be a list of Boolean words, and form its argument segments by
marking each word except the last as doubled. For every index $i<n$, segment $i$ is the
pair consisting of $w_i$ and the flag that is true exactly when $i+1<n$.
\end{lemma}

\begin{definition}[Oracle call specification]
\lean{Complexity.LogProg.CallSpec}
\label{Complexity.LogProg.CallSpec}
\leanok
\uses{Complexity.LogProg.Mode}
A call specification for a program with $m$ registers, $d$ deciders, and state set
$\Lambda$ consists of a decider indexed among the $d$ choices, a virtual-input mode, an
ordered list of register indices, and a state to enter after a positive answer and a
state to enter after a negative answer.
\end{definition}

\begin{definition}[Register program]
\lean{Complexity.LogProg.RProg}
\label{Complexity.LogProg.RProg}
\leanok
\uses{Complexity.LogProg.CallSpec, Turing.MultiTapeTM}
A register program with $m$ registers, $d$ deciders, and state set $\Lambda$ consists of
an $m$-work-tape Turing machine over the Boolean alphabet and a partial assignment of
oracle-call specifications to states.
\end{definition}

\begin{definition}[Virtual call segments]
\lean{Complexity.LogProg.callSegs}
\label{Complexity.LogProg.callSegs}
\leanok
\uses{Complexity.LogProg.CallSpec, Complexity.LogProg.Mode.seg0, Complexity.LogProg.Seg, Complexity.LogProg.argSegs}
Given a call specification, a Boolean input word $x$, and a word $W_i$ for each of its
$m$ registers, the call’s segment list consists of the initial segment selected by the
call mode and $x$, followed by the segments formed from the words of the registers
listed in the call specification, in their listed order.
\end{definition}

\begin{lemma}[Length of the Call Segmentation]
\lean{Complexity.LogProg.length_callSegs}
\label{Complexity.LogProg.length_callSegs}
\leanok
\uses{Complexity.LogProg.CallSpec, Complexity.LogProg.callSegs}
\difficulty{1}
For a call specification with argument list of length $q$, and for any input word and
assignment of words to registers, the resulting list of call segments has length $q+1$.
\end{lemma}

\begin{definition}[Program simulation states]
\lean{Complexity.LogProg.CSt}
\label{Complexity.LogProg.CSt}
\leanok
Given sets $\Lambda$ and $SD$ and a nonnegative integer $m$, define the program
simulation states to be the states of five forms: a program state labelled by an element
of $\Lambda$; a decider simulation state carrying a program label, a decider state, an
index in $\{0,\ldots,m\}$, two Boolean flags, and an optional Boolean result; two return
states, each carrying a program label, such an index, and two Boolean flags; and a
register return state carrying the same data as a return state together with a register
index in $\{0,\ldots,m-1\}$ and a Boolean scan flag.
\end{definition}

\begin{definition}[Signed virtual-tape step]
\lean{Complexity.LogProg.gstep}
\label{Complexity.LogProg.gstep}
\leanok
For Boolean flags
$\mathit{dbl},\mathit{last},\mathit{first},\mathit{par},\mathit{dir},\mathit{isChar}$
and a natural-number counter $s$, define a transition on signs as follows. Zero leaves
$(s,\mathit{par},\mathit{dir})$ unchanged and outputs zero. A positive sign, when
reading a character, changes the parity to true and outputs zero if $\mathit{dbl}$ is
true and parity is false; otherwise it sets parity to false, direction to true, and
outputs positive. For a positive sign not reading a character, if direction is true it
leaves the triple unchanged and outputs zero when $\mathit{last}$ is true, and otherwise
increments $s$, sets parity to true and direction to false, and outputs zero; if
direction is false it sets parity to false and direction to true and outputs positive. A
negative sign, when reading a character, sets parity to false and outputs zero if
$\mathit{dbl}$ is true and parity is true; otherwise it sets parity to true, direction
to false, and outputs negative. For a negative sign not reading a character, if
direction is true it sets parity to true, direction to false, and outputs negative; if
direction is false and $\mathit{first}$ is true it leaves the triple unchanged and
outputs zero; otherwise it decrements $s$, sets parity to false and direction to true,
and outputs zero.
\end{definition}

\begin{definition}[Clamping into a finite range]
\lean{Complexity.LogProg.toFin}
\label{Complexity.LogProg.toFin}
\leanok
For natural numbers $m$ and $s$, the associated element of $\mathrm{Fin}(m+1)$ is the
natural number $\min(s,m)$.
\end{definition}

\begin{definition}[Segment doubling flag]
\lean{Complexity.LogProg.segDbl}
\label{Complexity.LogProg.segDbl}
\leanok
\uses{Complexity.LogProg.CallSpec}
Given a call specification and a natural-number segment index $s$, define a Boolean flag
that is true exactly when either $s=0$ and the call mode is whole-input, or $s>0$ and
$s$ is less than the number of listed argument registers.
\end{definition}

\begin{definition}[Register assigned to a segment]
\lean{Complexity.LogProg.segReg}
\label{Complexity.LogProg.segReg}
\leanok
\uses{Complexity.LogProg.CallSpec}
Given a call specification with $m$ registers, a natural-number segment index $s$, and
$0<m$, assign to the segment the $(s-1)$-st listed register when that position exists,
and otherwise assign register $0$.
\end{definition}

\begin{definition}[Character-read flag]
\lean{Complexity.LogProg.isCharRead}
\label{Complexity.LogProg.isCharRead}
\leanok
\uses{Complexity.LogProg.Mode}
Given a mode, a natural-number segment index $s$, and an optional Boolean read symbol,
define a flag that is true exactly when the read symbol is true and $s=0$ in unary-first
mode; in every other case it is true exactly when a read symbol is present.
\end{definition}

\begin{definition}[Virtual tape symbol]
\lean{Complexity.LogProg.virtSym}
\label{Complexity.LogProg.virtSym}
\leanok
\uses{Complexity.LogProg.Mode, Complexity.LogProg.isCharRead}
Given a mode, a natural-number segment index $s$, a natural number $n$ of segments, a
direction flag, and an optional Boolean read symbol, define an optional virtual symbol
as follows. If the read counts as a character, return true when $s=0$ in unary-first
mode, and otherwise return the read symbol. If it does not count as a character, then in
the forward direction return no symbol when $s+1=n$, and false otherwise; in the reverse
direction return no symbol when $s=0$, and true otherwise.
\end{definition}

\begin{definition}[Register movement]
\lean{Complexity.LogProg.regMove}
\label{Complexity.LogProg.regMove}
\leanok
For $m\in\mathbb{N}$, a register index $r\in\{0,\ldots,m-1\}$, and a movement
$v\in\{-1,0,1\}$, define an action on the $m$ registers that writes nothing and moves
register $r$ by $v$, leaving every other register’s head stationary.
\end{definition}

\begin{definition}[Next return state]
\lean{Complexity.LogProg.nextRet}
\label{Complexity.LogProg.nextRet}
\leanok
\uses{Complexity.LogProg.CSt, Complexity.LogProg.CallSpec}
Given a call specification with argument list $A$, a program location $\ell$, a segment
index $s\in\{0,\ldots,m\}$, Boolean values $d$ and $b$, and a natural number $a$, define
the next control state to be the return state that restores register $a$ if $a<|A|$ and
$a<m$. Otherwise, it is the program state selected by $b$: the positive-answer state
when $b$ is true and the negative-answer state when $b$ is false.
\end{definition}

\begin{definition}[Idle decider action]
\lean{Complexity.LogProg.dIdle}
\label{Complexity.LogProg.dIdle}
\leanok
For a natural number $k_D$, define the action on each of $k_D$ decider work tapes to
write nothing and leave its head stationary.
\end{definition}

\begin{definition}[Compiled transition function]
\lean{Complexity.LogProg.ctr}
\label{Complexity.LogProg.ctr}
\leanok
\uses{Complexity.LogProg.CSt, Complexity.LogProg.RProg, Complexity.LogProg.dIdle, Complexity.LogProg.gstep, Complexity.LogProg.isCharRead, Complexity.LogProg.nextRet, Complexity.LogProg.regMove, Complexity.LogProg.segDbl, Complexity.LogProg.segReg, Complexity.LogProg.toFin, Complexity.LogProg.virtSym, Turing.Action, Turing.MultiTapeTM}
Let $P$ be a register program with $m$ registers and $d$ deciders, let $D$ be a Turing
machine with $k_D$ work tapes, and choose an initial state for each decider. Define a
transition function for control states consisting of ordinary program states,
call-simulation states, and return states, with $m+k_D$ work tapes. At an ordinary state
without a call, it performs the program’s transition on the first $m$ tapes and leaves
the remaining tapes idle; at a call, it initializes the selected decider simulation and
moves each listed argument-register head left. During simulation it runs the selected
decider on the encoded virtual input, applies the segment and head-movement update,
records the first emitted Boolean answer, and either continues simulation or begins the
return procedure if the decider halts. The return states rewind the input head and
restore the argument-register heads in order before continuing at the positive- or
negative-answer program state; transitions at states whose program location no longer
specifies a call halt.
\end{definition}

\begin{definition}[Compiled Turing machine]
\lean{Complexity.LogProg.compileTM}
\label{Complexity.LogProg.compileTM}
\leanok
\uses{Complexity.LogProg.CSt, Complexity.LogProg.RProg, Complexity.LogProg.ctr, Turing.MultiTapeTM}
Let $P$ be a register program with $m$ registers and $d$ deciders, let $\ell_0$ be its
initial program location, let $D$ be a Turing machine with $k_D$ work tapes, and choose
an initial state for each decider. Define the compiled Turing machine with $m+k_D$ work
tapes and control states for the compiled simulation: its initial state is the ordinary
program state at $\ell_0$, and its transition function is the compiled transition
function determined by $P$, $D$, and the chosen decider initial states.
\end{definition}

\begin{definition}[Embedding a machine configuration]
\lean{Complexity.LogProg.seam}
\label{Complexity.LogProg.seam}
\leanok
\uses{Complexity.LogProg.CSt, Turing.Cfg}
Given a configuration $c$ of a machine with $m$ work tapes on a Boolean input, define a
configuration of the compiled machine on the same input by mapping any present state to
the corresponding ordinary program state and retaining the input-head position and
output. The first $m$ work-tape contents and head positions are retained, while each of
the additional $k_D$ tapes is blank everywhere and has its head at position $0$.
\end{definition}

\begin{lemma}[Compiled machine initial configuration]
\lean{Complexity.LogProg.seam_init}
\label{Complexity.LogProg.seam_init}
\leanok
\uses{Complexity.LogProg.RProg, Complexity.LogProg.compileTM, Complexity.LogProg.seam, Turing.Cfg.init, Turing.MultiTapeTM, Turing.MultiTapeTM.initCfg}
\difficulty{4}
Let $P$ be a register program, $D$ a deterministic multi-tape machine, $l_0$ a program
state, and $q_0$ a map assigning a state of $D$ to each decider index. For every Boolean
input word $x$, the compiled machine's initial configuration on $x$ equals the seam
formed from the initial configuration of $P$ on $x$ with start state $l_0$.
\end{lemma}

\begin{lemma}[Compiled Step at a Program Node]
\lean{Complexity.LogProg.step_seam_prog}
\label{Complexity.LogProg.step_seam_prog}
\leanok
\uses{Complexity.LogProg.RProg, Complexity.LogProg.compileTM, Complexity.LogProg.ctr, Complexity.LogProg.dIdle, Complexity.LogProg.seam, Turing.Action.apply, Turing.Cfg, Turing.Cfg.inputSymbol, Turing.Cfg.workTapeSymbols, Turing.MultiTapeTM, Turing.MultiTapeTM.step}
\difficulty{5}
Let $P$ be a register program, $D$ a decider, and $q_0$ a choice of initial decider
state for each decider index. If a configuration $c$ of $P$ is at state $l$ and $l$ is
not a call node, then one step of the machine compiled from $P$ and $D$, started from
the configuration obtained by extending $c$ with idle decider tapes and marking its
state as the program state $l$, equals the extension of the configuration reached by one
step of $P$'s ordinary machine at $c$.
\end{lemma}

\begin{lemma}[Halting state under the seam construction]
\lean{Complexity.LogProg.seam_state_none}
\label{Complexity.LogProg.seam_state_none}
\leanok
\uses{Complexity.LogProg.seam, Turing.Cfg}
\difficulty{1}
Let $c$ be a configuration of an $m$-tape Turing machine on a fixed Boolean input, and
suppose its state is absent, which denotes that the machine has halted. Then the state
of the configuration obtained by the seam construction is also absent.
\end{lemma}
